\subsection*{Search restores without bypassing repentance\enspace\properrefs{Gosp.}}

\begin{fourstable}{Witness}{Searcher and lost}{Saving movement}{Precise sense}
St.~Cyril & Christ receives sinners who approach while accusers make mercy itself an accusation. & The physician uses nearness and table fellowship for healing, repentance, and joy. & Christ does not participate in their disorder; welcome has conversion as its end.\\
St.~Gregory & The Shepherd bears fallen humanity, and divine Wisdom lights and sweeps for the obscured image. & Incarnation searches; the shoulders bear; conscience is exposed; heaven rejoices. & The one and ninety-nine are theological figures, not demographic arithmetic.\\
St.~Ambrose & The shoulders become the Cross; the woman can signify the Church seeking the sovereign's coin. & Humanity lost in Adam is gathered in Christ, and ecclesial proclamation searches for the King's likeness. & Personal repentance belongs within common restoration; neither is erased by the other.\\
St.~Gregory of Nyssa & The lamp becomes illumined reason and the house the interior life. & Ascetical self-search uncovers a royal image obscured, not annihilated, by disorder. & Gregory of Nyssa reuses the parable for a subject other than Luke's, and his reading is not Gregory the Great's.\\
\end{fourstable}

\begin{sourcecard}{The two parables protect two truths}
The sheep prevents salvation from becoming autonomous introspection: the Shepherd comes from beyond the wanderer's resources and bears it home. The coin prevents \emph{lost} from meaning \emph{worthless}: the sovereign image grounds the search even while cleansing remains necessary. Both end in convoked joy, and Luke explicitly names repentance.
\end{sourcecard}

\subsection*{Recovered life enters common offering and joy\enspace\properrefs{Off., Sec., Comm., Postcomm.}}

\begin{threestable}{Liturgical moment}{Established movement}{Decisive reading}
\textbf{Offertory --- trust} & Psalm 9 turns those who know God's name toward trust, praise in Zion, and remembrance of the afflicted. & The found person does not become self-enclosed; received mercy opens an ecclesial voice toward other sufferers.\\
\textbf{Secret --- sanctification} & Personal \latin{Respice in me} becomes ecclesial \latin{Respice \ldots{} munera}; gifts are to be received for believers' salvation with enduring sanctification. & Irenaeus and Cyril of Jerusalem illuminate oblation and sanctified reception; neither comments on the Roman prayer.\\
\textbf{Communion --- shared joy} & Luke 15:10 places angelic joy over one repenting sinner on the communicants' lips. & Cyril and Gregory require the hearer to rejoice at another's restoration; sacramental joy remains joined to repentance.\\
\textbf{Postcommunion --- continuing effect} & Holy things already received are asked to vivify and prepare the cleansed for everlasting mercy. & Reception is neither a trophy nor a self-completing act; divine gift continues as purification, life, and preparation.\\
\end{threestable}

\begin{studybox}{Mercy, grace, dignity, and reception}
\begin{itemize}[topsep=0pt,itemsep=0.22em,parsep=0pt]
\item \textbf{Mercy is not indifference.} Luke's joy is joy over repentance under the just and patient Judge.
\item \textbf{Grace is not passivity.} The seeker acts first; the found person repents, resists, offers, and receives.
\item \textbf{Dignity is not self-restoration.} The image grounds the search; incarnation, Cross, illumination, and cleansing accomplish recovery.
\item \textbf{Reception is not automatic.} Secret and Postcommunion remain petitions.
\end{itemize}
\end{studybox}
